\section{Introduction}
\label{sec:introduction}

Skin-lesion classes can exhibit overlapping visual characteristics in
dermoscopic images, making automated multiclass classification difficult.
Deep convolutional neural networks have demonstrated strong representation
learning for skin-lesion analysis \cite{esteva2017dermatologist}, while
International Skin Imaging Collaboration (ISIC) challenges have supported
standardised development and evaluation on public dermoscopic collections
\cite{codella2019isic2018}.

Class imbalance remains a central problem. A classifier may achieve high
overall accuracy by recognising the majority class while performing poorly on
clinically important but under-represented categories. Evaluation should
therefore include balanced accuracy, macro-F1, and class-specific recall rather
than relying on accuracy alone.

Most multiclass systems make a single flat decision over the complete label
set. A conditional hierarchy instead divides the problem into broad and
fine-grained decisions. In the present setting, Stage~1 distinguishes
non-malignant from malignant lesions, and Stage~2 classifies predicted
malignant lesions as melanoma, basal cell carcinoma (BCC), or squamous cell
carcinoma (SCC).

This organisation introduces a routing dependency. A malignant lesion
incorrectly predicted as non-malignant is blocked before subtype
classification, while a non-malignant lesion incorrectly forwarded to
Stage~2 must receive a malignant label. Consequently, an apparently small
difference between flat and hierarchical endpoint scores may conceal a
substantial loss caused by the gate.

This paper studies that failure mode on a locked internal-test cohort derived
from ISIC 2019. The backbone-matched task-decomposition comparison uses
EfficientNet-B0 \cite{tan2019efficientnet}
for both a direct four-class model and a two-stage hierarchy.
The comparison does not isolate task decomposition alone, because the selected
Stage 2 classifier used class-balanced focal loss, whereas the flat
EfficientNet-B0 model used ordinary cross-entropy. A protocol-matched flat DenseNet-121
\cite{huang2017} is additionally evaluated on the same split and test cohort.
Only the convolutional backbone changes in this comparator; preprocessing,
cross-entropy training protocol, optimiser, scheduler, validation-based
checkpoint selection, and locked-test policy are preserved.

All three systems were evaluated on the same locked test cohort. Paired
bootstrap confidence intervals were used to quantify uncertainty in metric
differences for each pairwise comparison, while exact McNemar tests assessed
differences in paired correctness. The results are interpreted separately
for each metric rather than as evidence of uniform model superiority.

To examine the internal hierarchical failure mechanism, an oracle-routing
analysis retains the trained Stage~2 classifier but replaces the predicted
Stage~1 route with the ground-truth broad category. The resulting difference
estimates routing-associated loss while holding the downstream classifier
fixed. Route-specific rates further quantify malignant images blocked at
Stage~1, non-malignant images forwarded to Stage~2, and subtype errors among
correctly routed malignant images. Oracle routing is diagnostic only and is
not a deployable strategy.

A separate experiment evaluates Tis, T1, T2, T3, and T4 prediction on an
ISIC-derived melanoma T-category subset. It was not integrated into the
primary hierarchy and is reported only as standalone feasibility evidence
under severe label scarcity.

The contributions are:

\begin{enumerate}
    \item a same-cohort comparison of flat EfficientNet-B0, flat
    DenseNet-121, and a conditional EfficientNet-B0 hierarchy;

    \item paired statistical comparisons among the three systems using
    ground-truth-class-stratified bootstrap confidence intervals and exact
    McNemar tests; and

    \item an oracle-routing and route-specific decomposition that quantifies
    the contribution of Stage~1 gating errors to hierarchical endpoint
    performance.
\end{enumerate}

The remainder of this paper reviews related work, formalises the prediction
tasks, describes the experimental protocol, reports the results and
limitations, and concludes with priorities for future validation.
